% AY2019 力学 第2問 転記（問題のみ）
\section*{第2問〔II〕}

図2のように質量 $M$，慣性モーメント $Ma^2/2$，半径 $a$ の円板の周りに糸を巻きつけ，
糸の一端を固定して円板を落下させると，円板は反時計回りに回転し始めた。円板の重心 P は
$x$ 軸上に沿うものとする。ここで，鉛直下向きを $x$ 軸の正の向きとし，糸の張力を $T$，
重力加速度の大きさを $g$ とする。このとき，次の問い (1)〜(3) に答えよ。

\begin{enumerate}[label=(\arabic*)]
  \item 円板の重心 P の運動方程式を求めよ。
  \item 円板の角運動量の単位時間あたりの変化量の大きさを求めよ。
  \item 円板の重心 P の加速度の大きさを求めよ。
\end{enumerate}

\begin{center}
% 図2：天井に固定した糸を巻いた円板が落下。x軸=鉛直下向き正、重心P、半径a
\begin{tikzpicture}[>=stealth, scale=1.0]
  % 天井
  \fill[pattern=north east lines] (-0.4,3.0) rectangle (3.4,3.18);
  \draw (-0.4,3.0) -- (3.4,3.0);
  % 糸（右端で円板に接し鉛直に垂れる）
  \draw (2.0,3.0) -- (2.0,0);
  % 円板
  \draw (1.0,0) circle (1.0);
  \fill (1.0,0) circle (1.3pt); \node[left] at (0.9,0) {P};
  \draw[dashed] (1.0,0) -- (2.0,0); \node[above] at (1.5,0.0) {$a$};
  % x軸（鉛直下向き正）
  \draw[->] (1.0,1.7) -- (1.0,-2.2) node[below] {$x$};
\end{tikzpicture}

{\small 図 2}
\end{center}
